% 主旨：定义任务（鲁班七号镜像 1v1、首塔判胜）、350 维观测、6 分支条件动作、
%       训练协议与评价指标，为方法与实验提供统一记号。
\section{问题设定}
\label{sec:problem}

\subsection{对战任务}
\label{sec:task}

实验在腾讯开悟教育平台的王者荣耀 1v1 任务（版本 V63.0.9）上进行\cite{kaiwu63}。地图为单线的墨家机关道，两端是双方的泉水和水晶，水晶前各有一座外塔；双方英雄均为鲁班七号。外塔生命值为 \num{12000}，水晶为 \num{9000}。率先摧毁敌方外塔的一方获胜，满 \num{20000} 帧未分胜负记为超时。游戏每帧约 \qty{33.3}{\milli\second}，策略每 \num{3} 帧决策一次。地图存在战争迷雾，离开己方视野的敌方单位坐标被置为绝对值不小于 \num{100000} 的哨兵值。

\subsection{观测与动作}
\label{sec:obs-action}

平台特征处理把每帧原始状态编码为 \num{350} 维向量：双方英雄各 \num{67} 维（存活、等级、生命与法力比例、攻防属性、技能冷却、坐标与金币），双方各 \num{4} 个小兵槽位每槽 \num{18} 维，双方各 \num{2} 个防御建筑槽位每槽 \num{18} 维。V63.0.9 的任务 1 中小兵特征恒为零，与兵线相关的规则和奖励直接读取原始帧中的小兵坐标。

动作为 \num{6} 个离散分支 $a_t=(a_t^{(1)},\dots,a_t^{(6)})$，维数依次为 \num{12}、\num{16}、\num{16}、\num{16}、\num{16}、\num{9}。$a_t^{(1)}$ 为按钮（空动作、移动、普攻、三个技能、治疗、闪现、回城等）；$(a_t^{(2)},a_t^{(3)})$ 与 $(a_t^{(4)},a_t^{(5)})$ 分别把移动和施法方向量化为 $16\times16$ 栅格；$a_t^{(6)}$ 选择目标单位（无目标、敌方英雄、自身、\num{4} 个小兵槽、敌塔、野怪）。后五个分支是否生效以及目标分支的合法集合都取决于所选按钮。环境给出的 \num{184} 维合法动作向量包含按钮掩码、四个方向分支的掩码，以及按按钮分行的 $12\times9$ 目标掩码。

记 LSTM 隐状态为 $z_t$，分支 $k$ 的对数几率为 $\ell_\theta^{(k)}$，在前置分支 $a_t^{(<k)}$ 下的合法掩码为 $m_t^{(k)}(\cdot\mid a_t^{(<k)})\in\{0,1\}$。条件策略为
\begin{align}
  \pi_\theta^{(k)}(u\mid z_t,a_t^{(<k)},m_t)
  &=\frac{m_t^{(k)}(u\mid a_t^{(<k)})\exp\ell_\theta^{(k)}(u\mid z_t)}
         {\sum_{v}m_t^{(k)}(v\mid a_t^{(<k)})\exp\ell_\theta^{(k)}(v\mid z_t)},
  \label{eq:cond-policy-branch}\\
  \pi_\theta(a_t\mid z_t,m_t)
  &=\prod_{k=1}^{6}\pi_\theta^{(k)}(a_t^{(k)}\mid z_t,a_t^{(<k)},m_t).
  \label{eq:cond-policy}
\end{align}
训练时按 \cref{eq:cond-policy} 依次抽取全部六个分支，再按所选按钮决定哪些分支对环境生效，记生效指示为 $q_t^{(k)}\in\{0,1\}$；不生效的分支只在损失中权重为零（\cref{sec:base}）。评估时先对按钮取 $\argmax$，再在该按钮对应的合法目标中取 $\argmax$。

\subsection{训练协议与评价指标}
\label{sec:protocol}

平台以自对弈为主采样，红蓝方自动轮换；每 \num{10} 局自对弈后插入 \num{1} 局与内置规则智能体 \path{common_ai} 的评估局，评估局用确定性动作且不写入训练样本\cite{kaiwu63}。实验报告平台监控面板对 \path{common_ai} 评估局的统计：胜率、双方外塔剩余生命值（满值 \num{12000}）、场均击杀与死亡（K/D）、英雄伤抗比（对敌方英雄伤害与所受英雄伤害之比）和平均局长。监控把超时计为 \num{0.5} 胜，正式对战中超时双方均不得分，因此模型间对战另行报告胜、负、超时和异常的原始局数。
